\documentclass[10pt]{article}
\usepackage[T1]{fontenc}
\usepackage[utf8]{inputenc}
\usepackage[french]{babel}
\usepackage{amsmath,amssymb,mathrsfs,mathpazo}
\usepackage{geometry}
\geometry{a4paper,margin=2cm}
\usepackage{graphicx,array,tabularx,tabularray,colortbl,multirow,multicol,booktabs,enumitem,pifont,ulem,fancybox,diagbox,cancel}
\usepackage{pstricks,pstricks-add,pst-plot,pst-tree,pst-node,pst-text,pst-func,pst-math,pst-bspline,pst-3dplot,pst-eucl,pst-solides3d}
\let\pscalculate\undefined
\usepackage{pst-bezier}
\usepackage{tikz,pgfplots,tkz-tab}
\usetikzlibrary{arrows,arrows.meta,patterns,calc,positioning,trees,decorations.markings,shapes.geometric,shapes.misc,shapes.arrows,intersections,decorations.pathmorphing,decorations.pathreplacing,backgrounds,shadows,babel}
\pgfplotsset{compat=1.18}
\usepgfplotslibrary{fillbetween}
\usepackage[np]{numprint}
\usepackage{listings,makecell,pas-tableur,twoopt}
\let\barre\undefined
\usepackage[pstricks]{tablvar}
\definecolor{darkgreen}{RGB}{0,100,0}
\providecommand{\red}{\color{red}}
\providecommand{\blue}{\color{blue}}
\providecommand{\green}{\color{green}}
\providecommand{\black}{\color{black}}
\providecommand{\white}{\color{white}}
\providecommand{\gray}{\color{gray}}
\providecommand{\cyan}{\color{cyan}}
\providecommand{\magenta}{\color{magenta}}
\providecommand{\yellow}{\color{yellow}}
\providecommand{\violet}{\color{violet}}
\providecommand{\orange}{\color{orange}}
\providecommand{\euro}{\texteuro}
\providecommand{\EUR}{\texteuro}
\providecommand{\ieme}{\textsuperscript{e}}
\providecommand{\degres}{\ensuremath{^\circ}}
\providecommand{\degree}{\ensuremath{^\circ}}
\providecommand{\vv}[1]{\overrightarrow{#1}}
\providecommand{\Coord}[1]{#1}
\providecommand{\decosix}{}
\providecommand{\starredbullet}{\ensuremath{\bullet}}
\pagestyle{empty}
\setlength{\parindent}{0pt}
\usepackage{adjustbox}\DeclareUnicodeCharacter{2236}{:}\DeclareUnicodeCharacter{2190}{\ensuremath{\leftarrow}}
\providecommand{\R}{\mathbb{R}}
\providecommand{\N}{\mathbb{N}}
\providecommand{\Z}{\mathbb{Z}}
\providecommand{\Q}{\mathbb{Q}}
\providecommand{\D}{\mathbb{D}}
\providecommand{\C}{\mathbb{C}}
\providecommand{\e}{{\mathrm{e}}}
\providecommand{\vect}[1]{\overrightarrow{#1}}
\providecommand{\vectt}[1]{\overrightarrow{\mathrm{#1}}}
\providecommand{\vvt}[1]{\overrightarrow{\mathrm{#1}}}
\providecommand{\np}[2][]{{#2~#1}}
\providecommand{\vv}[1]{{\overrightarrow{#1}}}
\renewcommand{\d}{{\mathrm{d}}}
\providecommand{\degres}{{^\circ}}
\providecommand{\degre}{{^\circ}}
\providecommand{\blue}{\color{blue}}
\providecommand{\red}{\color{red}}
\providecommand{\enskip}{\quad}
\providecommand{\llbracket}{\lbrack\!\lbrack}
\providecommand{\rrbracket}{\rbrack\!\rbrack}
\providecommand{\euro}{{\text{€}}}
\providecommand{\no}{{\text{n°}}}
\providecommand{\interval}[2]{{[#1;#2]}}
\providecommand{\intervFO}[2]{{[#1;#2[}}
\providecommand{\dotfill}{\ldots}

\providecommand{\ds}{\displaystyle}
\providecommand{\textding}[1]{#1}
\providecommand{\Oij}{\(\mathrm{O};\vec{\imath},\vec{\jmath}\)}
\providecommand{\Oijk}{\(\mathrm{O};\vec{\imath},\vec{\jmath},\vec{k}\)}
\providecommand{\Ouv}{\(\mathrm{O};\vec{u},\vec{v}\)}
\providecommand{\intervFF}[2]{\ensuremath{[#1;#2]}}
\renewcommand{\euro}{\ifmmode\text{\texteuro}\else\texteuro\fi}
\DeclareUnicodeCharacter{00B0}{\ifmmode^{\circ}\else\textdegree\fi}
\DeclareUnicodeCharacter{20AC}{\ifmmode\text{\texteuro}\else\texteuro\fi}
% Énoncés : Éducation nationale. Transcription et corrigés : APMEP, Denis Vergès et contributeurs. Découpage et classement : CoopMaths, Sébastien Lozano, Éric Elter et contributeurs. Adaptation de la présentation : Maths Explorer.
% CC BY-SA : https://coopmaths.fr/www/about/
% https://coopmaths.fr/alea/static/e3c/2020/tex/e3c_2020_00_sujet23_3.tex
\begin{document}
\noindent\textit{APMEP / CoopMaths. CC BY-SA. Présentation adaptée par Maths Explorer.}\par\medskip
\medskip

\emph{Dans cet exercice, pour tout évènement A, on note $\overline{A}$ son évènement contraire, $p(A)$ sa probabilité et, si $B$ est un évènement de probabilité non nulle, $p_B(A)$ la probabilité conditionnelle de $A$ sachant $B$.}

\medskip

Une entreprise a fabriqué en un mois \np{1500} chaudières, dont $900$~chaudières à
cheminée et $600$~chaudières à ventouse.

On a constaté, dans ce lot, que :

\begin{itemize}
\item 1\,\% des chaudières à cheminées ont un défaut
\item 6\,\% des chaudières à ventouses ont un défaut.
\end{itemize}

On prélève au hasard le numéro de série d'une chaudière de la production de ce
mois.

On considère les évènements suivants :

\begin{itemize}
\item $C$ : \og Le numéro de série est celui d'une chaudière à cheminée \fg
\item $V$ : \og Le numéro de série est celui d'une chaudière à ventouse \fg
\item $D$ : \og Le numéro de série est celui d'une chaudière défectueuse \fg
\end{itemize}

\medskip

\begin{enumerate}
\item Recopier et compléter sur la copie le tableau à double entrée suivant :

\medskip

\begin{tabularx}{\linewidth}{|m{3.3315cm}|*{3}{X|}}
\cline{2-4}
\multicolumn{1}{c|}{}&
nombre de chaudières à cheminée & nombre de chaudières à ventouse &
 \centering Total\tabularnewline\hline
nombre de chaudières
défectueuses
&&&\\\hline
nombre de chaudières non défectueuses
&&&\\\hline
 Total &\centering 900& \centering 600& \centering 1500\tabularnewline
\hline
\end{tabularx}

\item Recopier et compléter l'arbre pondéré suivant :

\begin{center}
\psset{nodesepA=0pt,nodesepB=3pt,treesep=0.75,labelsep=0.1pt}
\pstree[treemode=R]{\TR{}}
{\pstree{\TR{$C$~}\taput{$\np{0.6}$}}
	{
	\TR{$D$}\taput{$\dots$}
	\TR{$\overline{D}$}\tbput{$\dots$}
	}
\pstree{\TR{$V$~}\tbput{$\dots$}}
	{\TR{$D$}\taput{$\dots$}
	\TR{$\overline{D}$}\tbput{$\dots$}
	}
}
\end{center}

\item Calculer la probabilité que le numéro de série soit celui d'une chaudière
défectueuse.
\item Déterminer $P_D(V)$. Interpréter le résultat dans le contexte de l'exercice.
\item Les évènements $D$ et $V$ sont-ils indépendants ?
\end{enumerate}

\vspace{0,5cm}

\end{document}
